\section{Coordinate kernels and incidence geometries}\label{coordinate-kernels-and-incidence-geometries}

An incidence system has a finite type set, objects carrying types, and a
symmetric incidence relation; distinct objects of one type are not
incident. A flag is a set of pairwise incident objects. A chamber has
one object of each type. A geometry is an incidence system in which
every flag extends to a chamber. The residue of a flag consists of the
objects of the remaining types incident with that flag, with induced
incidence. The geometry is residually connected when the incidence graph
of every residue of rank at least two, including the empty-flag residue,
is connected. It is thin when every panel, a flag missing just one type,
lies in exactly two chambers. A regular hypertope is a thin, residually
connected geometry admitting a chamber-transitive type-preserving
automorphism group.

For a marked group indexed by \(I\), its maximal parabolics are
\(P_i=G_{I\setminus\{i\}}\). The associated typed coset incidence system
has objects \(P_i g\), and incidence means nonempty intersection of
cosets. All actions and cosets below use this right-coset convention.

For subgroups \(P_i\leq G\), typed coset Helly means that every pairwise
intersecting family of right cosets \(P_i g_i\), with distinct indices,
has a common element. This property is separate from IP. It supplies a
common representative for an incident flag, whereas IP determines the
stabilizer of a flag already represented in this way.

\textbf{Proposition UA.Geometry (coset recognition and faithful chamber
action).} Suppose the marking of \(G\) consists of nonidentity
involutions, has IP, and has typed coset Helly for its maximal
parabolics. The coset incidence system is then a thin, residually
connected geometry, and \(G\) is its full type-preserving automorphism
group, acting regularly on chambers. More generally, the full
type-preserving automorphism group of any thin, residually connected
geometry of finite rank acts freely on chambers. A chamber-transitive
subgroup is consequently the full group.

\emph{Proof.} In the coset system, a common coset point extends every
flag to a chamber by choosing the coset of each missing type through
that point. Right translation sends every such flag to its base flag.
For a base flag of type set \(S\), IP gives stabilizer
\(G_{I\setminus S}\). In particular the full chamber stabilizer is
\(1\), while a panel stabilizer is a group generated by one nonidentity
involution and has order two. There are therefore exactly two chambers
through each panel.

Consider a residue with remaining type set \(T\), where \(|T|\geq2\).
Its stabilizer is \(G_T\), and its maximal object stabilizers are
\(G_{T\setminus\{i\}}\), \(i\in T\). Every marker indexed by \(T\) lies
in at least one of these proper parabolics. A word in those subgroups
supplies a walk between their translated base objects in the residue
incidence graph; all objects are translates of base objects and all base
objects lie in a common flag. The residue is therefore connected.

For the general assertion, first observe that every residually connected
finite-rank geometry has a connected chamber graph. This follows by
induction on rank: chambers through a fixed object are connected through
galleries in its residue, and a path in the incidence graph, extending
each consecutive incident pair to a chamber, connects arbitrary chambers
through these subgraphs. The rank-one case holds because all chambers
share the empty panel. In a thin geometry an automorphism fixing a
chamber fixes its unique neighbor across each panel. It therefore fixes
every chamber along every gallery, and hence every object. The chamber
action is free. If a subgroup is chamber-transitive, any full
automorphism agrees on a base chamber with one subgroup element; their
quotient fixes that chamber and is the identity. Apply this to the coset
geometry to obtain the stated full automorphism group. \(\square\)

\subsection{Helly transfer through coordinate
kernels}\label{helly-transfer-through-coordinate-kernels}

Write right cosets as \(P g\), with the subgroup on the left. Let
\(q:H\to B\) be an epimorphism of groups generated by corresponding
marked families \((h_i)_{i\in I}\) and \((b_i)_{i\in I}\), where \(I\)
is finite. For \(J\subseteq I\), write \(H_J=\langle h_j:j\in J\rangle\)
and \(B_J=\langle b_j:j\in J\rangle\). In particular, \(q(H_J)=B_J\) for
every J.

Fix any set \(\mathcal J\) of subsets of I as the permitted types. The
maximal-parabolic case is \(\mathcal J=\{I\setminus\{i\}:i\in I\}\).

\textbf{Theorem UA.SK.Helly (coordinate-kernel transfer).} Suppose B has
IP. Suppose \(N=\ker q\) is an internal direct product
\(\prod_{\alpha\in A}N_\alpha\), with a finite factor set A, and for
every \(J\in\mathcal J\), \[
 N\cap H_J=\prod_{\alpha\in A_J}N_\alpha
 \qquad\text{for some }A_J\subseteq A.
 \tag{H1}
\] If the family \((B_J)_{J\in\mathcal J}\) has typed coset Helly, then
so does \((H_J)_{J\in\mathcal J}\).

No commutativity within the factors, finiteness of H or B, or
faithfulness of an action of B on N is assumed. The required
intersection-image fact is \(q(H_{J\cap K})=B_J\cap B_K\); it follows
here from the marked map and IP of B. IP of H is not an additional
assumption for this transfer theorem.

\emph{Proof.} Consider pairwise intersecting cosets \(H_{J_i}g_i\).
Their images \(B_{J_i}q(g_i)\) intersect pairwise. Choose a common base
element b by the assumed Helly property, and choose u in H with
\(q(u)=b\). Each set \[
 F_i=H_{J_i}g_i\cap q^{-1}(b)
\] is nonempty: because \(bq(g_i)^{-1}\in B_{J_i}\), a preimage of that
element in \(H_{J_i}\), multiplied on the right by \(g_i\), belongs to
\(F_i\).

We first show that \(F_i\cap F_j\ne\varnothing\). Choose
\(x\in H_{J_i}g_i\cap H_{J_j}g_j\) and put z=q(x). Both b and z belong
to each of the two base cosets, so \[
 bz^{-1}\in B_{J_i}\cap B_{J_j}=B_{J_i\cap J_j}.
\] The marked map is surjective on this intersection parabolic. Choose
\(h\in H_{J_i\cap J_j}\) with \(q(h)=bz^{-1}\). Since h belongs to both
\(H_{J_i}\) and \(H_{J_j}\), left multiplication by h preserves both
right cosets. Thus hx belongs to their intersection and has image b.
This proves the assertion without assuming that q carries an arbitrary
intersection onto the intersection of its images.

Choose \(x_i\in F_i\) and put \(k_i=x_i u^{-1}\in N\). If \(y\in F_i\),
then \(yx_i^{-1}\in N\cap H_{J_i}\); the converse also holds. Hence \[
 F_i u^{-1}=(N\cap H_{J_i})k_i.
 \tag{H2}
\] Under the direct product coordinates of N, (H1) makes (H2)
unrestricted in coordinates \(\alpha\in A_{J_i}\), and fixes coordinate
\(\alpha\notin A_{J_i}\) to \((k_i)_\alpha\). Pairwise intersections
imply that any two constraints fixing the same coordinate fix it to the
same value. Choose that value in every constrained coordinate and the
identity in every unconstrained coordinate. The resulting element
belongs to all sets in (H2). Multiplication on the right by u supplies a
common element of the original cosets. This proves the theorem.
\(\square\)

\textbf{Corollary UA.SK.Geometry.} In the supported-kernel realization
UA.SK, assume the base marking gives a regular hypertope and every
marked generator in every attached group is an involution. Then the
resulting marking also gives a regular hypertope, and its full
type-preserving automorphism group is the constructed group.

\emph{Proof.} UA.SK gives IP and the exact profile \[
 N\cap H_J=\prod_{S\subseteq J}N_S.
\] The base maximal-parabolic cosets have typed Helly. Theorem
UA.SK.Helly applies with factors indexed by the supports S. Each
constructed generator is a tuple whose base coordinate and every
nontrivial attached coordinate are involutions; hence its square is the
identity. Its base image is nonidentity, so its order is exactly two.
UA.Geometry now gives thinness, residual connectedness, chamber
regularity and equality with the full type-preserving automorphism
group. In particular, this conclusion uses typed Helly in addition to
IP. \(\square\)

A sufficient alternative to local involutivity is that every support has
size at least two: each singleton local parabolic is then proper, and
injectivity into its nonidentity involutory base image forces the local
marker to have order two. This is a sufficient condition, not an
equivalent reformulation. Without either condition the geometric
corollary fails. With a singleton type set and support, the quotient
\(C_4\to C_2\) is injective on its only proper parabolic, the trivial
subgroup. UA.SK then constructs \(C_4\), whose distinguished generator
has order four. In rank three, adjoining two independent \(C_2\) markers
gives \(C_4\times C_2^2\) with the same defective singleton attachment
and a four-element panel.
